% No preamble. Prefix gdd-. Fixed endpoint, actual determinant densities.
\section{Gauge-invariant diagonal determinant densities}
\label{gdd-section}

Fix a composite endpoint \(K=\ell^k\), a preceding homogeneous
depth \(M\), and the actual finite periodic construction. The
endpoint is fixed throughout this section. Constants may depend
on \(K\); they will be uniform in the permitted volume and in
the preceding depth where the upstream classical bounds are.
Choose the fixed sufficiently large scalar \(a>0\) from the
coarse determinant estimate of Notebook 32.

At the actual stationary field \(U_*(W)\), use its left-relative
objects
\[
 C=f'(0),\quad s'=C^T\lambda,\quad
 \mathcal L=s''-\lambda\cdot f'',\quad
 G_b=D_b^TD_b,\quad
 Q_b=\mathcal L+D_bD_b^T+aC^TC,\quad
 P_c=q_{\mathrm{eff}}+aI.
\]
The exact finite identity already established is
\begin{equation}
 D_{M,K}(W)=\mathcal N_0\left[
       \log\det G_b-\frac12\log\det Q_b
                           +\frac12\log\det P_c\right].
 \label{gdd-exact-total}
\end{equation}
Here \(\mathcal N_0\) subtracts the value at \(W=1\).
The matrices act, respectively, on noncenter scalar sites,
current bonds, and output bonds, including their color coordinates.

\subsection{Background covariance of the actual matrices}
\label{gdd-covariance-section}

For a fixed current-site gauge \(u\), let \(O_E\) denote
the block-diagonal adjoint action at each current bond's
initial site, \(O_S\) the corresponding action at noncenter
scalar sites, and \(O_C\) its action at output bond initial
centers. These matrices are orthogonal for real gauges.
They are complex orthogonal on the holomorphic continuation.
The transformation of a left-relative variation is linear:
\[
 (e^{ix}U)^u=e^{iO_Ex}U^u.
\]
The actual ordinary average is equivariant, so its relative
target logarithm transforms by \(O_C\). Local gauge invariance
of the two actual classical actions consequently gives
\begin{align}
 C(U^u)&=O_C C(U)O_E^T,&
 \lambda(W^{u|_{\rm centers}})&=O_C\lambda(W),\notag\\
 \mathcal L(U^u)&=O_E\mathcal L(U)O_E^T,&
 D_b(U^u)&=O_E D_b(U)O_S^T. \label{gdd-transformations}
\end{align}
For the Lagrangian identity, differentiate
\[
 s_{U^u}(O_Ex)=s_U(x),\qquad
 f_{U^u}(O_Ex)=O_Cf_U(x).
\]
The second equality and the multiplier transformation retain
the entire \(\lambda\cdot f''\) term. No unconstrained
stationarity is assumed here.

The based scalar domain is unchanged by these transformations:
zero values at all \(K\)-centers remain zero. The effective
relative Hessian transforms by \(O_C\), since its target
relative variable also transforms linearly. Thus
\begin{equation}
 G_b(U^u)=O_SG_b(U)O_S^T,\quad
 Q_b(U^u)=O_EQ_b(U)O_E^T,\quad
 P_c(W^{u|_{\rm centers}})=O_CP_c(W)O_C^T.
 \label{gdd-matrix-covariance}
\end{equation}
These identities concern the fixed scalar completion \(aI\).
They do not assert such covariance for an arbitrary
background-dependent completion.

\subsection{Fixed-endpoint locality of the matrix logarithms}
\label{gdd-log-locality-section}

We need locality of a diagonal logarithm, not merely locality
of the trace of its derivatives. The required fixed-endpoint
spectral and kernel inputs follow from the preceding modules,
as follows.

At zero background, \(G_b\) has a positive gap independent
of volume when \(K\) is fixed. Indeed, join every site to
its block center by a path of length at most \(c_dK\).
For a scalar field vanishing at centers, Cauchy--Schwarz
along these paths and summation in a block give
\begin{equation}
 \|\phi\|_2^2\le c_d K^{d+1}\|D_b\phi\|_2^2.
 \label{gdd-pinned-Poincare}
\end{equation}
The crude exponent is sufficient here. This is compatible
with the soft-mode upper bound in Notebook 32 and supplies
no gap uniform as \(K\) grows.

The flat constrained contour precision has its previously
proved positive gap at fixed \(K\), uniformly in preceding
depth and volume. On \(\ker C\), decompose a vector into
its component in \(\operatorname{ran}D_b\) and its orthogonal
complement. The form \(\mathcal L\) annihilates the first.
For a horizontal vector \(v\in\ker C\), replace it by
\(v+D_bTv\) to make its contour derivative zero.
Its \(\mathcal L\) value is unchanged, and the orthogonal
projection of this replacement back to the horizontal
space is \(v\). The contour precision bound therefore
controls \(\|v\|_2^2\). Its passage from retained
coordinates to current bonds uses the fixed-\(K\)
bounded insertion \(K_0\); this only changes the constant.
On vertical vectors \(D_b\phi\), the divergence penalty
has the lower bound
\[
 \|D_b^TD_b\phi\|_2^2
       \ge \lambda_{\min}(G_b)\|D_b\phi\|_2^2.
\]
Thus \(\mathcal L+D_bD_b^T\) has a uniform fixed-\(K\)
gap on \(\ker C\).

The actual adjusted minimum-section derivative
\[
 H_b=H-D_bG_b^{-1}D_b^TH
\]
has \(CH_b=I\), \(D_b^TH_b=0\), and a bounded operator norm
at fixed \(K\). The latter follows from the established
minimum-response bound and the orthogonal projection
in this formula. Write \(x=\xi+H_by\), \(C\xi=0\).
The exact effective Schur identity gives
\[
 x^TQ_bx
  =\xi^T(\mathcal L+D_bD_b^T)\xi
                         +y^T(q_{\mathrm{eff}}+aI)y.
\]
Since \(q_{\mathrm{eff}}+aI\) has a fixed positive gap,
this proves a flat gap for \(Q_b\), uniform in volume
and preceding depth at fixed \(K\). Their upper spectral
bounds and exponentially weighted matrix bounds follow
from the actual strict classical families, the finite-range
constraint at fixed \(K\), and the local gauge generator.
\(P_c\) already has the stronger uniform bounds of Notebook 32.

Compose these matrices with the actual stationary section.
Its strict support kernels, together with those of the
classical actions and local charts, give source-rooted
matrix derivative bounds through order two. A sufficiently
small common complex output ball preserves the flat
gaps in operator norm. This ball may depend on \(K\),
but not on volume or preceding depth.

Here is an elementary logarithm argument making the
locality consequence precise. For any one of these
matrix families \(P(w)\), choose a fixed scalar \(s\)
larger than its flat upper spectral bound and put
\[
 Z(w)=I-P(w)/s.
\]
At zero \(\|Z(0)\|_{2\to2}<1\), uniformly over the
volumes under consideration. The exponentially weighted
kernel bound implies that conjugating \(Z\) by a small
exponential distance weight changes its operator norm
by an arbitrarily small amount. First use bounded
truncations of the distance weight; their uniform
Lipschitz bound gives the same estimate, and then pass
to kernel bounds. No unbounded similarity is assumed
to act on all of \(\ell^2\). Shrink the background
ball as well. The conjugated operators then have norm
at most one fixed \(r_K<1\). Therefore
\begin{equation}
 \log P(w)=(\log s)I-\sum_{n\ge1}\frac{Z(w)^n}{n}
 \label{gdd-log-series}
\end{equation}
converges with these exponential weights.
Termwise differentiation through order two is permitted:
there are at most \(n\) or \(n^2\) insertions, absorbed
by the geometric factor \(r_K^{n-2}\).
The actual source-rooted bounds attach every background
source to the matrix-product chain. Exponential weights
then connect the diagonal matrix coordinate and all
sources. Conversion from weighted \(\ell^2\) to row
sums spends a smaller exponent and a finite geometric
sum. The internal fiber per output block is finite
at fixed \(K\).

In particular, if a diagonal color trace is assigned
to a fixed output block \(x\), its quadratic kernel
\(k_x(b,c)\) obeys, for some \(\kappa_K>0\),
\begin{equation}
 \sum_{b,c}e^{\kappa_K\tau(x,b,c)}|k_x(b,c)|\le C_K.
 \label{gdd-rooted-kernel}
\end{equation}
Here the sum includes directions and colors, \(\tau\)
is a connecting-tree distance in output units, and all
finite multiplicities of current coordinates are included.
Periodic bounds use shortest lifted distances. No extensive
logarithm is taken to an infinite-volume limit.

\subsection{A density with its own gauge identity}
\label{gdd-density-section}

Assign every noncenter current site to its unique centered
\(K\)-block. Assign each positive current bond to the block
containing its initial site, and each output bond to its
initial output site. Define
\begin{align}
 \rho_x(w)=\mathcal N_0\bigg[
 &\sum_{\substack{z\ {\rm noncenter}\\z\ {\rm in\ block}\ x}}
                   \operatorname{tr}_{\mathfrak g}(\log G_b)_{zz}
 -\frac12\sum_{\substack{e\ {\rm current}\\e_-\ {\rm in\ block}\ x}}
                   \operatorname{tr}_{\mathfrak g}(\log Q_b)_{ee}
 \notag\\
 &+\frac12\sum_{\substack{c\ {\rm output}\\c_-=x}}
                   \operatorname{tr}_{\mathfrak g}(\log P_c)_{cc}
 \bigg].
 \label{gdd-density}
\end{align}
All matrices in this formula are evaluated at \(U_*(e^{iw})\)
or its output \(e^{iw}\), as appropriate. On the finite
regulator,
\begin{equation}
                         D_{M,K}(e^{iw})=\sum_x\rho_x(w)
 \label{gdd-density-total}
\end{equation}
exactly.

Functional calculus respects the similarities in
\eqref{gdd-matrix-covariance}. A diagonal color trace is
unchanged by its site's orthogonal conjugation. Moreover,
the actual minimum at gauge-transformed output is gauge
equivalent to the transformed old minimum, on the established
local germ. To see this, transform the constraint and action,
then use the proved uniqueness of the local constrained
critical orbit and its selected contour representative.
The comparison need only hold on a finite real germ;
analyticity supplies the local identities used below.
It follows that \emph{each} \(\rho_x\), and each of its
three displayed groups of terms, has its own local
output gauge identity. This conclusion uses more than
gauge invariance of their sum.

The block assignment is translation covariant. Local
convergence of the matrices and their derivatives, together
with \eqref{gdd-log-series} and its uniform geometric tail,
passes these densities and their quadratic kernels to
any jointly compatible infinite system at fixed endpoint.
There they are translation covariant, satisfy
\eqref{gdd-rooted-kernel}, and retain the local gauge identity.
No uniqueness of different structural limits is needed.

\subsection{An affine-density formula for the total Maxwell moment}
\label{gdd-moment-section}

\begin{lemma}[Gauge-invariant rooted density and total moment]
\label{gdd-moment-lemma}
Let \(\rho_x\) be translation-covariant analytic scalar
densities on an infinite lattice, with exponentially
rooted quadratic kernels as in \eqref{gdd-rooted-kernel}.
Assume every \(\rho_x\) is invariant under local gauge
changes on its germ, and under global semisimple color
conjugation. Fix one color and distinct directions
\(\mu,\nu\). For the formal affine component
\[
                  A_\mu(z)=t z_\nu,\qquad A_\sigma=0
                    \quad(\sigma\ne\mu),
\]
the quadratic density coefficient is
\begin{equation}
 [t^2]\rho_0(A)
       =-\frac14\sum_r r_\nu^2 K_{\mu\mu}(0,r),
 \label{gdd-affine-moment}
\end{equation}
where \(K=\sum_x D^2\rho_x(0)\) is the total Hessian
kernel in that color. Both sides are absolutely convergent
kernel expressions. The left side does not mean evaluation
of an analytic map at an unbounded nonzero field.
\end{lemma}

\begin{proof}
Global semisimple color invariance kills each density's
linear derivative. Differentiating its local gauge identity
at the origin therefore shows that its Hessian annihilates
each compactly supported infinitesimal gradient in either
slot. The exponential kernel bound extends this identity
to polynomial gauge parameters by absolute summation.
In particular it annihilates a constant one-form in any
one color and direction: use the linear scalar gauge
parameter with that gradient.

Write \(k_0(b,c)\) for the selected same-component Hessian.
Its row and column sums are zero. Consequently
\[
 \frac12\sum_{b,c} k_0(b,c)b_\nu c_\nu
   =-\frac14\sum_{b,c}k_0(b,c)(c_\nu-b_\nu)^2.
\]
Every interchange is allowed by exponential rooted
summability. Translation covariance gives
\[
 K(0,r)=\sum_x k_x(0,r)=\sum_x k_0(-x,r-x).
\]
The change of indices \(b=-x,\ c=r-x\) turns the last
second moment into the right side of
\eqref{gdd-affine-moment}. The same calculation with
the translated affine field shows that the coefficient
is independent of the selected density root.
\end{proof}

For the actual correction, the established Maxwell
normalization is
\(\beta(D)=-\tfrac12\sum_r r_\nu^2K_{\mu\mu}(0,r)\),
in one normalized color. Thus the constructed density gives
\begin{equation}
                  \boxed{\ \beta(D_{M,K})
                            =2[t^2]\rho_0(tz_\nu e_\mu)\ }.
 \label{gdd-beta-density}
\end{equation}
The density construction applies separately to the
three groups of terms in \eqref{gdd-density}, with the
corresponding second-moment extraction on each.
When the total has the established spatial symmetries,
this is its Maxwell coefficient.

The individual gauge identity is essential. For example,
the Maxwell action may be written either as its curvature
square density or, after summation by parts, as a density
of the form \(\tfrac12 A\cdot d^*dA\). On an affine
constant-curvature field the latter vanishes pointwise,
while the former does not. Their totals agree on compact
fields. The diagonal densities above avoid this ambiguity
because their own gauge identities imply the row and
column cancellations used in the proof.

\subsection{Scope}
\label{gdd-scope}

The result identifies a specific, locally gauge-invariant
trace density with the total Maxwell moment. It does not
estimate that moment uniformly as \(K\) grows. The gap,
analytic radius and decay exponent used for \(G_b,Q_b\)
are permitted to deteriorate with \(K\). No positive
coefficient or comparison with a massive reference
follows from the density identity alone.
